%% kapa: exam=1 round=25 year=2014 period=1 form=A subject=civil_law subject_ko=민법 seq=19 number=19
%% source: Q-Net 감정평가사 기출문제 (cst00309) artlSeq=5101207 file=extracted/1008/5101207/2014년 제25회 감정평가사 1차시험 기출문제/1교시(A형).hwp
%% conversion: kordoc 4.24.51 -> tools/gen1.py (auto); review=proofread-sonnet; flags=none
\begin{question}{19}
사기․강박에 의한 의사표시에 관한 설명으로 옳지 않은 것은? (다툼이 있으면 판례에 의함)
\begin{choices}
  \item 상대방 있는 의사표시에 관하여 제3자가 사기를 행한 경우에 표의자는 상대방이 그 사실을 안 경우에 한하여 그 의사표시를 취소할 수 있다.
  \item 피기망자에게 손해를 가할 의사는 사기에 의한 의사표시의 성립요건이 아니다.
  \item 강박에 의한 의사표시가 스스로 의사결정을 할 수 있는 여지가 전혀 없는 상태에서 의사표시의 외형만 있는 것에 불과한 경우에 그 의사표시는 효력이 없다.
  \item 상대방이 불법적인 해악의 고지 없이 각서에 서명ㆍ날인할 것을 강력히 요구하는 것만으로도 강박이 된다.
  \item 교환계약의 일방 당사자가 자기 소유의 목적물의 가액을 시가보다 높게 허위로 고지하였더라도 특별한 사정이 없는 한, 기망행위에 해당하지 않는다.
\end{choices}
\end{question}
\answer{1,4}
\answernote{복수 정답(정답표 원문: 1, 4)}
